% GENERATED FILE. Edit canonical lessons, then run npm run generate:books.
% canonical-source-hash: fnv1a64:40ca7469cdb5e6c0
% canonical-lessons: ML-C295-kukkar, ML-C295-miksi, ML-C295-ammi, ML-C295-ural, ML-C295-ulakka

\chapter{Pressure cooker, Mixer-grinder, Grinding stone, Mortar, Pestle}
\label{ch:ml-pressure-cooker-mixer-grinder-grinding-stone-mortar-pestle}
\hlchaptermodality{295}

\begin{chapteropening}
\textbf{By the end of this chapter:} \emph{I can say a pressure cooker, a mixer-grinder, a grinding stone, a mortar and a pestle.}

Complete the last lesson of chapter 295: I can say a pressure cooker, a mixer-grinder, a grinding stone, a mortar and a pestle.
\end{chapteropening}

\section[\texorpdfstring{\ml{കുക്കർ} (kukkar)}{കുക്കർ (kukkar)}]{\ml{കുക്കർ} (kukkar) — a pressure cooker}

\label{lesson:ML-C295-kukkar}



\begin{quote}
Before the new one: say the Malayalam for a chain, then the Malayalam for a stick.
\end{quote}

\subsection*{You'll want to know: \ml{കുക്കർ}}
\textbf{\ml{കുക്കർ}} — \emph{kukkar} — \textquotedblleft{}a pressure cooker\textquotedblright{}.

\begin{grammarlens}[title={What you've built}]
One more word for this chapter.
\end{grammarlens}

\subsection*{Your turn}
\begin{itemize}
\raggedright
  \item \emph{Say it:} \emph{kukkar}
  \item \emph{Say it:} \emph{kukkar}, once more
  \item \emph{Say it:} say \emph{kukkar}
  \item \emph{Recall:} say the Malayalam for a chain, then the Malayalam for a stick, then say \emph{kukkar} again
\end{itemize}

\begin{tcolorbox}[breakable,colback=teal!4,colframe=teal!35!black,title={Before you move on}]
What does \ml{കുക്കർ} mean? (\textquotedblleft{}A pressure cooker\textquotedblright{}.) Say it once more.
\end{tcolorbox}
\section[\texorpdfstring{\ml{മിക്സി} (miksi)}{മിക്സി (miksi)}]{\ml{മിക്സി} (miksi) — a mixer-grinder}

\label{lesson:ML-C295-miksi}



\begin{quote}
Before the new one: say the Malayalam for a stick, then the Malayalam for a pressure cooker.
\end{quote}

\subsection*{You'll want to know: \ml{മിക്സി}}
\textbf{\ml{മിക്സി}} — \emph{miksi} — \textquotedblleft{}a mixer-grinder\textquotedblright{}.

\begin{grammarlens}[title={What you've built}]
One more word for this chapter.
\end{grammarlens}

\subsection*{Your turn}
\begin{itemize}
\raggedright
  \item \emph{Say it:} \emph{miksi}
  \item \emph{Say it:} \emph{miksi}, once more
  \item \emph{Say it:} say \emph{miksi}
  \item \emph{Recall:} say the Malayalam for a stick, then the Malayalam for a pressure cooker, then say \emph{miksi} again
\end{itemize}

\begin{tcolorbox}[breakable,colback=teal!4,colframe=teal!35!black,title={Before you move on}]
What does \ml{മിക്സി} mean? (\textquotedblleft{}A mixer-grinder\textquotedblright{}.) Say it once more.
\end{tcolorbox}
\section[\texorpdfstring{\ml{അമ്മി} (ammi)}{അമ്മി (ammi)}]{\ml{അമ്മി} (ammi) — a grinding stone}

\label{lesson:ML-C295-ammi}



\begin{quote}
Before the new one: say the Malayalam for a pressure cooker, then the Malayalam for a mixer-grinder.
\end{quote}

\subsection*{You'll want to know: \ml{അമ്മി}}
\textbf{\ml{അമ്മി}} — \emph{ammi} — \textquotedblleft{}a grinding stone\textquotedblright{}.

\begin{grammarlens}[title={What you've built}]
One more word for this chapter.
\end{grammarlens}

\subsection*{Your turn}
\begin{itemize}
\raggedright
  \item \emph{Say it:} \emph{ammi}
  \item \emph{Say it:} \emph{ammi}, once more
  \item \emph{Say it:} say \emph{ammi}
  \item \emph{Recall:} say the Malayalam for a pressure cooker, then the Malayalam for a mixer-grinder, then say \emph{ammi} again
\end{itemize}

\begin{tcolorbox}[breakable,colback=teal!4,colframe=teal!35!black,title={Before you move on}]
What does \ml{അമ്മി} mean? (\textquotedblleft{}A grinding stone\textquotedblright{}.) Say it once more.
\end{tcolorbox}
\section[\texorpdfstring{\ml{ഉരൽ} (ural)}{ഉരൽ (ural)}]{\ml{ഉരൽ} (ural) — a mortar (for pounding)}

\label{lesson:ML-C295-ural}



\begin{quote}
Before the new one: say the Malayalam for a mixer-grinder, then the Malayalam for a grinding stone.
\end{quote}

\subsection*{You'll want to know: \ml{ഉരൽ}}
\textbf{\ml{ഉരൽ}} — \emph{ural} — \textquotedblleft{}a mortar (for pounding)\textquotedblright{}.

\begin{grammarlens}[title={What you've built}]
One more word for this chapter.
\end{grammarlens}

\subsection*{Your turn}
\begin{itemize}
\raggedright
  \item \emph{Say it:} \emph{ural}
  \item \emph{Say it:} \emph{ural}, once more
  \item \emph{Say it:} say \emph{ural}
  \item \emph{Recall:} say the Malayalam for a mixer-grinder, then the Malayalam for a grinding stone, then say \emph{ural} again
\end{itemize}

\begin{tcolorbox}[breakable,colback=teal!4,colframe=teal!35!black,title={Before you move on}]
What does \ml{ഉരൽ} mean? (\textquotedblleft{}A mortar (for pounding)\textquotedblright{}.) Say it once more.
\end{tcolorbox}
\section[\texorpdfstring{\ml{ഉലക്ക} (ulakka)}{ഉലക്ക (ulakka)}]{\ml{ഉലക്ക} (ulakka) — a pestle}

\label{lesson:ML-C295-ulakka}



\begin{quote}
Before the new one: say the Malayalam for a grinding stone, then the Malayalam for a mortar.
\end{quote}

\subsection*{You'll want to know: \ml{ഉലക്ക}}
\textbf{\ml{ഉലക്ക}} — \emph{ulakka} — \textquotedblleft{}a pestle\textquotedblright{}.

\begin{grammarlens}[title={What you've built}]
One more word for this chapter.
\end{grammarlens}

\subsection*{Your turn}
\begin{itemize}
\raggedright
  \item \emph{Say it:} \emph{ulakka}
  \item \emph{Say it:} \emph{ulakka}, once more
  \item \emph{Say it:} say \emph{ulakka}
  \item \emph{Recall:} say the Malayalam for a grinding stone, then the Malayalam for a mortar, then say \emph{ulakka} again
\end{itemize}

\begin{tcolorbox}[breakable,colback=teal!4,colframe=teal!35!black,title={Before you move on}]
What does \ml{ഉലക്ക} mean? (\textquotedblleft{}A pestle\textquotedblright{}.) Say it once more.
\end{tcolorbox}
